\documentclass[11pt]{article}

% ============================================================
% Response to Reviewer -- Round 22
% Companion to: revised_paper_v18.tex
% "Mortgage Lock-In and the Federal Reserve's Quantitative
%  Tightening Shortfall"
% Sixteen-item report across three tiers.
% Section references are to revised_paper_v18.tex; every numeric
% value is drawn from a committed run artifact or from the code
% path named inline.
% ============================================================

\usepackage[margin=1in]{geometry}
\usepackage{amsmath}
\usepackage{amssymb}
\usepackage{enumitem}
\usepackage{parskip}
\usepackage[hidelinks]{hyperref}

\newcommand{\mlabel}[1]{\texttt{#1}}
\newcommand{\run}[1]{\texttt{#1}}

\title{\vspace{-2em}Response to Reviewer\\[0.3em]
\large Manuscript: \emph{Mortgage Lock-In and the Federal Reserve's
Quantitative Tightening Shortfall}\\[0.2em]
\normalsize Sixteen-item report: Danish counterfactual, fit statistics,
floor form, and estimator scope}
\author{}
\date{26 July 2026}

\begin{document}
\sloppy
\maketitle
\vspace{-2.5em}

Dear Editor and Reviewer,

Thank you for the report. It is the most useful I have had, and I want to be
straightforward about two things before the point-by-point.

First, several items were written against a draft that a prior revision had
already changed. I say so not to deflect but because the reader deserves to know
which of my responses are new work and which are pointers to work already in the
manuscript; I have labelled every item accordingly, and where an item is already
met I quote the sentence that meets it so you can judge for yourself rather than
take my word.

Second, and more important: the items that did land were not marginal. Your
functional-form point (item 3) turned out to have a second half neither of us
had seen, your fit-statistic point (item 2) was right about a defect I had been
reporting for three rounds, and your instinct about the Danish leg (item 1) was
correct on a channel different from the one you named. Chasing them produced
five new pre-committed runs and, along the way, four defects in the manuscript
that no reviewer had raised and my own verification suite could not see. Those
are itemised in \S5, because a report that surfaces errors it did not itself
identify has done its job better than a report that only lists its own hits.

Numbers below are quoted from committed artifacts. New runs follow the project's
spec-before-run convention: a specification commit precedes each result commit,
and each carries parity gates that replay committed values bit-exactly before
anything new is reported.

\section{Item 1 --- the Danish counterfactual}

\textbf{Status: the constructive fix was already in the manuscript; one prong
you raised was live, and it was worse than you argued.}

\subsection{1(a) The floor violation and the U.S.-intercept transplant}

This is implemented, and by exactly the construction you prescribe. The
production leg (\run{danish\_us\_intercept}) keeps the identified flatness but
anchors the level at the U.S. hazard evaluated at zero rate gap, with the
involuntary floor retained: \texttt{competing\_risks.py} passes a zeroed rate gap
into \texttt{prepay\_hazard}, so the floor is applied on the Danish leg exactly as
on the U.S.\ leg. It lands at 5.61\% mean Danish CPR --- your own predicted
``roughly 6\% CPR'' --- and turns the gap positive: $+\$61.2$ billion at the
in-sample calibration point and $+\$28.2$ billion at the off-window calibration
the headline marginal uses.

The Danish-level leg survives only as a labelled bracketing case, and your
diagnosis of it is correct: it prepays at 3.39\%, below the 4\% floor both
frameworks impose. That inconsistency is now disclosed wherever the
$-\$99.9$ billion figure appears, including the Introduction and the Conclusion,
which previously carried the number without the caveat.

\subsection{1(b) ``No Danish mechanism'': the numbers are right, the inference
is not}

Your arithmetic checks exactly --- I verified it from the two parquets rather
than from the tables. The rule-only Danish leg and the $\beta_1 = 0$ null agree
to $0.000107$ percentage points of mean CPR and to $\$0.004$ billion of trapped
balance.

But that identity is forced, and it is forced by the fix you ask for in 1(a). The
gap enters the hazard only through $(-\beta_1) \times \text{gap}$; setting the gap
to zero and setting $\beta_1$ to zero are the same operation on that term. A
gap-flat hazard at the U.S.\ zero-gap intercept \emph{is} the null leg with the
floor retained. So the coincidence is not evidence that the mechanism is absent;
it is what a correct transplant must produce. The manuscript states this
directly: the two are ``one object measured on two accounting legs.''

\subsection{1(c) The discount-rate table --- you were right, and about a larger
channel than you named}

You read the exactly-zero deltas as evidence that market-value repurchase is
absent rather than as immunity to a pricing critique. That reading is correct,
and pursuing it found two things.

The smaller one is that the diagnostic had executed on the \emph{wrong leg}.
\texttt{danish\_discount\_bound.py} contains no call to
\texttt{set\_danish\_moving\_anchor}, so it ran at the module default
(\texttt{dk\_level}); its artifact records a mean Danish CPR of 3.39\%, the
bracketing anchor's speed, not the production leg's 5.61\%. The table caption
nevertheless described it as the production leg and extended the result to the
production gap. Caption and text now name the leg that ran, and state that the
production anchor's PV-independence is structural rather than run-verified.

The larger one is the accounting. The Danish leg credits every retired balance at
\emph{par}: the dynamic-balance loop books roll-off as balance times drain with
no price term anywhere in it. Netting the roughly \$199 billion scheduled
component from the leg's \$669.3 billion of Danish roll-off leaves on the order
of \$470 billion credited at par, against a proxy discount of 32--34\% at the
representative 3.0\%/6.8\% state. A market-value credit would reduce Danish cash
receipts by more than the $+\$61.2$ billion gap itself, and in the direction that
shrinks it. I have not restated the gap under such a credit, because the pipeline
does not distinguish a cash haircut from a balance adjustment; I record it as the
largest un-priced channel in the counterfactual and retract the claim that the
pricing critique attaches to the superseded variant alone.

\subsection{1(d)--1(e) The sweep, and the Berger citation}

The sign-fragility framing you object to is gone: the figure and the text now say
that the anchor, not the refinance channel, decides the sign. Two further
corrections came out of re-reading that figure, both reported in \S5.

On your falsifier --- whether Berger et al.\ estimate 3.2\% per year as a
general-equilibrium prediction for U.S.\ households or as a Danish descriptive
level --- I checked the source directly. It is the latter: the January 2026 draft
uses 3.2\% as ``the unconditional average moving rate among Danish FRM
borrowers,'' a calibration target in its \S4.6. The manuscript's slope-versus-level
distinction is therefore the correct reading of that paper, and it now carries
locators so you can check it in one step.

\section{Item 2 --- recovery percentage as the primary fit statistic}

\textbf{Status: implemented, with one correction to the argument.}

You are right that recovery is an affine transform of cumulative error with a
large additive constant, and that this flatters every estimator. \mlabel{tab:bases}
now carries a terminal cumulative runoff error column, the note states the
identity explicitly, and errors are given as shares of the \$652.8 billion of
runoff actually realized. Your arithmetic reproduces exactly.

One correction. Nothing changes sign under the transform: recovery
$= 1 + \text{error}/764.7$ is monotone. What flips is the \emph{referent} --- an
over-prediction of trapped liquidity is an under-prediction of runoff --- and the
denominator. I have written it that way because the sharper claim is available
without the stronger one: Path B's in-sample $+\$53.8$ billion is simultaneously
an 8.2\% under-prediction of realized runoff, and the headline off-window leg,
which reads as a worse fit at 91.3\%, in fact carries the \emph{smallest} error in
the table on the standalone scorer, $+\$2.8$ billion or 0.4\%.

Your asymmetry point about WSHOMCB rescaling is a correct reading of the code and
now underwrites a different result; see item 6.

\section{Item 3 --- the floor's functional form}

\textbf{Status: the level-versus-form test was already in the manuscript; the
prong you identified as still open was open, and it refuted a claim I was making
implicitly.}

The additive competing-risks form had been run at the production floor and at the
three off-window anchors, giving a form-conditional identified range of $+3.9$ to
$+13.1$ points. Your specific charge --- that the concave-transform check had
been executed inside the censoring regime it needed to escape --- was correct:
\run{concave\_marginal} ran under the hard-maximum form only.

Running the joint cell (\run{concave\_additive\_marginal}) leaves the range intact:
the concave $\times$ additive marginal is $+9.25$ points at the production floor and
$+9.22$ at the headline anchor, inside the hull at both. But it refutes something
the manuscript had been asserting without measurement. \mlabel{tab:uncertainty}
listed the form adjustment and the transform adjustment side by side, which
presumes they add. They do not: the interaction is $-1.47$ points at the headline
anchor, past the $\pm 1$-point convention used throughout. The mechanism is the
censoring at issue --- the concave transform costs $-0.52$ points under the max
form, which absorbs most of it into the floor, against $-1.99$ under the additive
form, which never censors. The table now reports the joint cell and the
interaction rather than implying additivity.

I have also disclosed what the additive form costs in level terms, which the
manuscript had priced only at the production floor: at the off-window anchors the
additive central leg recovers 55.9\% of the benchmark against a 44.7\% null, and
the hull's upper endpoint comes from a cell recovering 61.2\%. The
form-conditional range is therefore a range over specifications of sharply
different aggregate fit, and its upper end is not an equally credentialed member.
That cuts against the additive form rather than for it, and it belongs in the
open.

\section{Item 4 --- the ABM; item 5 --- Path A; item 16 --- scale}

\textbf{Item 4: partly already met, one prong newly closed.} The abstract already
reported the fifty-seed mean rather than the seed-42 draw. Your remaining charge
--- that the cross-design threshold test was a single draw from an estimator with
material seed dispersion --- was correct, and \run{cross\_design\_seeds} now runs
fifty seeds across both randomization layers. The verdict survives and
strengthens: mean 60.2\%, standard deviation 3.3 points, 95\% band 55.1--66.0\%,
the committed draw at the 48th percentile of its own distribution, and all fifty
seeds classifying as undercutting. The band is reported wherever the point
appears.

That sweep also separates the two variants on the manuscript's own admissibility
test, in the direction that was not convenient for me: the recalibrated variant's
deep-gap floor is inside the observed 4--5\% band on all fifty seeds, and the
frozen variant's is outside it on all fifty. Applied consistently, my own test
disqualifies the leg I had been reporting as corroborating.

\textbf{Item 5: your premise does not hold, and I want to be exact about where.}
The abstract's 97.9\% is Path~B's in-sample recovery, not Path~A's; the abstract
contains no Path~A figure. But your underlying concern was sound at one site:
Section~I carried Path~A's 112.4\% describing it as corroboration, and a
bias-respecting sign test (\run{patha\_sign\_test}) does not reject
$\beta_g \leq 0$ under any construction. The corroboration language is retracted
there and in the interpretation section, and the diagnostic table now labels each
row's specification --- it had been interleaving spec v4 and spec v3 rows without
saying so.

\textbf{Item 16: correct, and unmet until now.} The scale check exercised the
household mechanic in isolation, without the DTI wall, the freeze trait,
curtailment, multi-vintage weighting or the settlement kernel, and ran 3.9 points
hot against the production ABM. The claim is now scoped at all five sites to what
was measured. A production-specification rerun is reported separately; its first
execution halted at a pre-committed parity tier on upstream data revision, and I
have recorded the halt rather than relaxing the tolerance.

\section{Defects this report surfaced indirectly}

These were found while chasing the items above. None was raised in the report and
none was visible to the verification suite, which pins numeric literals and
therefore cannot see a false statement about a magnitude, an omitted summary
statistic, or a caption describing the wrong object.

\begin{enumerate}[leftmargin=1.4em]
\item \textbf{A false claim in the sweep figure's caption.} It asserted that both
anchors stay ``an order of magnitude below the superseded \$925.5 billion
extrapolation across the entire sweep.'' The body text five lines earlier gives
$+\$1{,}038.9$ billion at the ceiling, and both artifacts agree; both anchors
\emph{exceed} \$925.5 billion there. Corrected, with both ceilings quoted.

\item \textbf{Negative trapped balances in the same sweep.} At refinance
contributions of 12\% and above the Danish leg's trapped balance goes to
$-\$14.4$, $-\$158.0$ and $-\$290.0$ billion. That is not an economically
meaningful quantity, and the upper half of the axis is now marked directional only.

\item \textbf{Selective reporting in the Path A bootstrap.} The manuscript
reported the median, the interquartile range, the interval and the fraction of
draws above the benchmark, and omitted the mean: \$598.9 billion, or 78.3\% of
the benchmark --- below the paper's own mechanical null. Now reported, with the
reason.

\item \textbf{An ablation that had never been run.} The burnout ablation was
quoted at four sites with no artifact, script or gate behind it, and its figures
reproduce the \emph{minimum} over 100 replicates of an unrelated origination-time
permutation experiment. Executed properly (\run{burnout\_ablation}), the true
value is $-\$6.58$ billion against the claimed $-\$6.6$ billion: the number was
right by coincidence. The run also showed the ablation halves at the headline
calibration, where three interpretive sites had been quoting the in-sample figure.
\end{enumerate}

\section{Items 6--15, and what I have not done}

Item~6's global validation statement has been moved into Section~I. The
loan-level bootstrap you ask for existed, but you were right that it was the
wrong scheme: it resampled within strata and held the 130 strata fixed, so
between-stratum variance was never drawn. Drawing the strata
(\run{bootstrap\_pathb\_cluster}) widens the marginal's interval by roughly a
factor of 37, to $[+4.63, +6.92]$ points. The floor-read layer remains the wider
of the two, so the headline is unaffected --- but ``the single stratified draw
contributes essentially no uncertainty'' was a statement about the scheme, not
the estimator, and it is retracted.

Item~14's third pre-commitment gap was the sharpest thing in the report. The
admissibility precondition excluding the externally anchored variant does not
appear in that run's specification commit, which instead states that the variant
``is evaluated against the SAME pre-registered bands'' and pre-registers the floor
diagnostic as reported ``whichever direction it moves''; the frozen artifact
classifies the result as undercutting. The manuscript now states the
pre-committed classification first and labels the qualification as post hoc, and
the verification suite reads the artifact's classification so a green suite can no
longer pass over it.

Items~7, 8, 10, 11, 12(b), 12(c) and 15 were met in the manuscript already, and I
quote the operative sentences in the point-by-point appendix rather than here.
Item~13's reduction is correct and is now stated explicitly: the no-trade-off
claim reduces to the marginal being bounded by the excess of the null hazard over
the floor, and to zero where the floor binds --- a property of the production form
that lifts under the additive one.

Two items I have \emph{not} implemented, with reasons. Item~9's abstract remains
longer than you would like: cutting it further requires deleting a disclosed
result, and each candidate is a disclosure a prior round added deliberately. It
now stands at 424 words. And item~12(a)'s Ginnie overlay is promoted into
Table~1 but not made the headline level, because \$27.8 billion of the raw
adjustment is a component common to any observed-series overlay; I have labelled
it an accounting-correction variant instead, and I am glad to go further if you
think that understates it.

\vspace{1em}
\noindent Thank you again. The report cost me more runs than any previous one and
the manuscript is materially more honest for it.

\end{document}
